\documentclass{INSAToulouse}

\renewcommand{\familydefault}{\sfdefault} % Police Arial-like

\usepackage{listings, xcolor, array, pgfgantt, lmodern, multirow, tabularx, colortbl, adjustbox}
\usetikzlibrary{positioning, fit, backgrounds, arrows.meta, decorations.pathreplacing, shapes.multipart, calc}
\definecolor{mainred}{HTML}{FF4C4C}


% ── À remplir ────────────────────────────────────────────────────────────────
\titre{PROJET DE FIN D'ÉTUDES}
\imagecouverture{example-image}

%\enabletlssec % à décommenter pour les TLS-SEC

\firstcouverture{
    \parbox{\textwidth}{
        \sffamily
        \textbf{BigD}\\
        Informatique et Réseaux\\
        Département GEI\\
        Spécialité TLS-SEC\\
        Promotion 59
    }
}

\secondcouverture{
    \parbox{\textwidth}{
        \begin{custombox}{9cm}
            \sffamily
            \textbf{DÉVELOPPEMENT D'UN SYSTEME DE CHIFFREMENT BASÉ SUR LE PUR HASARD}
            \vspace{1.5em}\\
            %\textbf{Lieu du PFE}\\
            Bullshit corp\\
            15 Avenue du papillon\\
            31400 Toulouse
            \vspace{1em}\\
            \textbf{Tuteur d'alternance}\\
            Tic
            \vspace{1em}\\
            \textbf{Correspondant pédagogique INSA}\\
            Tac
            \vspace{1em}\\
            \textbf{Soutenance du PFE le 3/09/2026}
        \end{custombox}
    }
}
% ─────────────────────────────────────────────────────────────────────────────

\usepackage{biblatex}
\addbibresource{bibliography.bib}

% ── Glossaire ─────────────────────────────────────────────────────────────────
% Ajouter vos définitions ici, puis utiliser \gls{cle} dans le texte
\newglossaryentry{pki}{name={PKI}, description={Public Key Infrastructure}}
% ─────────────────────────────────────────────────────────────────────────────
\makeglossaries  % doit être dans le préambule, avant \begin{document}

\begin{document}

    \makecover          % page de couverture (fond PDF INSA)
    \makepagegarde      % page de garde intérieure

    \pchapter{Résumé}
    blabla

    \clearpage

    \pchapter{Remerciements}
    blabla
    
    \clearpage

    \pchapter{Note sur l'utilisation d'outils externes}
    j'ai utilisé l'IA, blabla
    \clearpage

    \makebackmatter     % liste des figures + liste des tableaux
    \makeglossaire      % lexique / glossaire

    \clearpage
    
    \makesommaire       % table des matières + passage en numérotation arabe

    %%%%%%%%%%%%%%%%%%%%%%%
    %% INTRODUCTION :    %%
    %%%%%%%%%%%%%%%%%%%%%%%
    \pchapter{Introduction}
    blabla
    oui y'a une \gls{pki}

    %%%%%%%%%%%%%%%%%%%%%%%
    %% CHAPTER 1    :    %%
    %%%%%%%%%%%%%%%%%%%%%%%
    \chapter{J'explique le hasard}
    \section{Explication physique}
    \section{Faut être très précis}
    \section{Explication bancale}
    jador l'ANSSI et surtout l'ii901\cite{ii901}
    
    %%%%%%%%%%%%%%%%%%%%%%
    %% CONCLUSIONS :    %%
    %%%%%%%%%%%%%%%%%%%%%%
    \pchapter{Conclusion}


    % appendix etc
    \appendix
    \chapter{voici une annexe}
    blabla

    
    % bibliographie
    \clearpage
    \printbibliography[heading=bibintoc, title={Bibliographie}]
    
    \makebackcover      % dernière page (fond PDF INSA)

\end{document}
